\title{DeepSeekMath: Pushing the Limits of Mathematical Reasoning in Open Language Models}

\begin{abstract}
Mathematical reasoning remains a major obstacle for language models owing to its inherent complexity and structured demands. In this study, we present DeepSeekMath~7B, a model derived by further pre-training DeepSeek-Coder-Base-v1.5 7B on 120B math-centric tokens extracted from Common Crawl, supplemented with both natural language and code datasets. DeepSeekMath~7B attains a notable 51.7\% on the competition-grade MATH benchmark, operating independently of auxiliary toolkits or ensemble voting strategies and attaining performance close to that of Gemini-Ultra and GPT-4. Utilizing 64-sample self-consistency, DeepSeekMath~7B reaches a score of 60.9\% on MATH. Two principal strategies underlie the model’s strong mathematical reasoning: First, an advanced data selection mechanism capitalizes on the wealth of openly accessible web content; Second, Group Relative Policy Optimization (GRPO)—an adaptation of Proximal Policy Optimization (PPO)—is proposed to boost reasoning proficiency while also improving PPO’s memory efficiency.
\end{abstract}

\section{Introduction}

Large language models (LLMs) have transformed methodologies for mathematical reasoning within artificial intelligence, leading to notable progress on quantitative reasoning benchmarks \citep{MATH} as well as geometry reasoning evaluations \citep{trinh2024solving}. In addition, these systems have become valuable tools for aiding humans with challenging mathematical tasks \citep{tao}. Nonetheless, advanced systems like GPT-4 \citep{gpt4} and Gemini-Ultra \citep{gemini} remain closed to the public, while open-source alternatives that are available exhibit significantly lower performance.

In this paper, we present DeepSeekMath, a specialized language model tailored for mathematical tasks, which demonstrates substantial improvements over existing open-source models and attains performance nearing that of GPT-4 when evaluated against academic standards. This advancement is supported by the construction of the DeepSeekMath Corpus, an extensive and meticulously curated pre-training dataset containing 120B mathematical tokens. Our data collection process involves extracting content from the Common Crawl (CC) via a fastText-driven classifier \citep{joulin2016fasttext}. Initially, this classifier is trained with positive samples drawn from OpenWebMath \citep{openwebmath} and a heterogeneous mix of unrelated web documents as negatives. Using the trained classifier, we subsequently extract further positive candidates from the CC; these are subjected to additional filtering and validation through human annotation. The augmented data subsequently informs the retraining of the classifier for higher precision. Assessment of the resulting corpus confirms its quality: our foundational model, DeepSeekMath-Base 7B, achieves 64.2\% accuracy on GSM8K \citep{gsm8k} and 36.2\% on the competitive MATH benchmark \citep{MATH}, surpassing the results reported for Minerva 540B \citep{minerva}. Moreover, owing to the multilingual composition of the DeepSeekMath Corpus, notable gains are observed on Chinese mathematical evaluations \citep{wei2023cmath,agieval}. We propose that our approach to mathematical data curation provides a valuable reference point for future efforts within the field, and we anticipate further advancements in subsequent research.

DeepSeekMath-Base is constructed by initializing with DeepSeek-Coder-Base-v1.5 7B \citep{deepseek-coder}, based on our findings that commencing from a code-pretrained model yields superior results relative to a general-purpose LLM. Additionally, our experiments show that mathematical pretraining boosts the model’s performance not only on mathematical tasks but also on general knowledge and reasoning benchmarks such as MMLU \citep{mmlu} and BBH \citep{bbh}, suggesting an overall enhancement in reasoning abilities.

Subsequent to pre-training, DeepSeekMath-Base is subjected to mathematical instruction fine-tuning, employing chain-of-thought \citep{cot}, program-of-thought \citep{pot,pal}, and tool-augmented reasoning data \citep{tora}. This fine-tuning process yields DeepSeekMath-Instruct 7B, which surpasses all other 7B parameter models and performs on par with much larger, 70B-parameter, open-source models tuned with instructions.

Additionally, we present Group Relative Policy Optimization (GRPO), a novel variant of Proximal Policy Optimization (PPO) \citep{schulman2017proximal} within the reinforcement learning (RL) landscape. GRPO eliminates the need for a critic by utilizing group scores to estimate the baseline, which substantially lowers the computational requirements for training. Leveraging only a portion of the English instruction tuning dataset, GRPO achieves marked performance gains relative to the strong DeepSeekMath-Instruct baseline, demonstrating notable improvements on both in-domain datasets (GSM8K: 82.9\% $\rightarrow$ 88.2\%, MATH: 46.8\% $\rightarrow$ 51.7\%) and out-of-domain benchmarks (for instance, CMATH: 84.6\% $\rightarrow$ 88.8\%) during RL fine-tuning. Furthermore, we establish a unified framework for interpreting various approaches, such as Rejection Sampling Fine-Tuning (RFT) \citep{yuan2023scaling}, Direct Preference Optimization (DPO) \citep{dpo}, PPO, and GRPO. This unified perspective reveals that these algorithms can be understood as either direct or streamlined forms of RL. Our comprehensive experimental analysis—involving comparisons such as online vs. offline RL, outcome-based vs. process-based supervision, and single-step vs. multi-step RL—sheds light on the core drivers behind this framework. Lastly, we provide an explanation for the observed performance improvements from RL on instruction-tuned models and outline future research avenues for developing more robust RL methodologies in light of our unified perspective.

\subsection{Contributions}
Our contribution includes scalable math pre-training, along with the exploration and analysis of reinforcement learning.

\textbf{Math Pre-Training at Scale}

\begin{itemize}[topsep=0pt]
    \item In this work, we demonstrate that the openly available Common Crawl dataset harbors significant resources pertinent to mathematics. Through the use of a carefully crafted data selection process, we have assembled the DeepSeekMath Corpus—a large-scale, high-quality dataset comprising 120B tokens sourced from web pages specifically curated for mathematical relevance. This corpus is nearly 7 times larger than the math web pages utilized by Minerva \citep{minerva} and exceeds the recently released OpenWebMath \citep{openwebmath} by a factor of 9.

\item DeepSeekMath-Base 7B, our foundational pre-trained model, attains a performance level similar to that of Minerva 540B \citep{minerva}. This result suggests that mathematical reasoning proficiency is not determined solely by the model's parameter count. Notably, even models of modest size can demonstrate robust mathematical reasoning skills when trained on carefully curated high-quality data.

\item We present our observations from experiments focused on mathematical training. Pretraining models on code before exposing them to mathematical tasks leads to enhanced performance in solving mathematical problems, regardless of whether external tools are employed. This provides some evidence addressing the long-debated issue: \textit{does exposure to code improve reasoning skills?} Our results suggest that, for mathematical reasoning specifically, it indeed has a positive effect.

    \item Although training on arXiv papers is common, especially in many math-related papers, it brings no notable improvements on all mathematical benchmarks adopted in this paper.
\end{itemize}

\textbf{Exploration and Analysis of Reinforcement Learning}

\begin{itemize}[topsep=0pt]
    \item We present Group Relative Policy Optimization (GRPO), a reinforcement learning approach that is both resource-efficient and effective. Distinct from Proximal Policy Optimization (PPO), GRPO eliminates the need for a critic model by utilizing group scores to estimate the baseline, thereby substantially lowering the computational overhead during training.
    \item Our results show that GRPO provides considerable improvements to our instruction-tuned model, DeepSeekMath-Instruct, leveraging only instruction-tuning datasets. Moreover, the reinforcement learning phase leads to noticeable gains in out-of-domain generalization.
    \item We establish a comprehensive framework for interpreting various reinforcement learning strategies, including RFT, DPO, PPO, and GRPO. Through a range of thorough empirical studies—spanning topics such as online versus offline training, outcome versus process-based supervision, and single-turn versus iterative reinforcement learning—we examine the critical components within this framework.
    \item Leveraging our unified perspective, we investigate the underlying mechanisms driving the success of reinforcement learning, and outline multiple promising pathways for enhancing reinforcement learning in large language models.
\end{itemize}

\subsection{Summary of Evaluations and Metrics}

\begin{itemize}[topsep=0pt]
    \item \textbf{English and Chinese Mathematical Reasoning}: Our models are rigorously evaluated across a variety of English and Chinese benchmark datasets, spanning mathematical tasks from elementary school through university level. For English, the assessment draws on datasets such as GSM8K \citep{gsm8k}, MATH \citep{MATH}, SAT \citep{llemma}, OCW Courses \citep{minerva}, and MMLU-STEM \citep{mmlu}. The Chinese evaluation utilizes benchmarks like MGSM-zh \citep{mgsm}, CMATH \citep{wei2023cmath}, Gaokao-MathCloze \citep{agieval}, and Gaokao-MathQA \citep{agieval}. The evaluation examines not only the capacity of our models to produce self-contained textual solutions in the absence of external tools, but also their proficiency in tackling mathematical problems using Python-based solutions.

On English evaluation tasks, DeepSeekMath-Base demonstrates performance on par with the proprietary Minerva 540B \citep{minerva}, and outperforms all publicly available base models (such as Mistral 7B \citep{mistral} and Llemma-34B \citep{llemma}), irrespective of their exposure to math pre-training, frequently by a considerable margin. Particularly, DeepSeekMath-Base exhibits a clear advantage on Chinese tasks, which can be attributed to our approach of not limiting the math pre-training corpus to English, as done in prior studies \citep{minerva,llemma}, and instead incorporating high-quality math data in various languages. After applying mathematical instruction tuning and reinforcement learning, the enhanced models—DeepSeekMath-Instruct and DeepSeekMath-RL—achieve strong results, surpassing the 50\% accuracy threshold on the competition-level MATH benchmark, marking a first for open-source models.

\item \textbf{Formal Mathematics}: We assess DeepSeekMath-Base on the informal-to-formal theorem proving benchmark introduced by \citep{dsp_proof}, utilizing the miniF2F dataset \citep{minif2f} and selecting Isabelle \citep{isabelle} as the proof assistant. The results indicate that DeepSeekMath-Base achieves robust performance in few-shot autoformalization scenarios.
\item \textbf{Natural Language Understanding, Reasoning, and Code}: To holistically evaluate general capabilities in comprehension, reasoning, and programming, we test DeepSeekMath-Base across several standard benchmarks: Massive Multitask Language Understanding (MMLU) \citep{mmlu}, which features 57 diverse multiple-choice problems; BIG-Bench Hard (BBH) \citep{bbh}, comprised of 23 tasks designed to require advanced multi-step reasoning; as well as HumanEval \citep{codex} and MBPP \citep{mbpp}, both widely adopted for assessing code generation models. Pre-training with mathematical data enhances the model’s performance not only in language understanding but also in complex reasoning tasks.

\section{Math Pre-Training}

\subsection{Data Collection and Decontamination} This section details the methodology employed to build the DeepSeekMath Corpus utilizing data from Common Crawl. As illustrated in Figure \ref{fig:data_collect}, we describe an iterative workflow that methodically accumulates a large-scale mathematical text corpus from Common Crawl, beginning with an initial high-quality seed dataset focused on mathematics. It should be emphasized that this methodology can be similarly extended to other specialized fields, such as programming.

Initially, we select OpenWebMath \citep{openwebmath}—a repository of curated mathematical web documents—as our starting corpus. Using this dataset, we construct and train a fastText model \citep{joulin2016fasttext} to retrieve additional web pages similar in mathematical content to OpenWebMath. For the training process, we randomly draw 500,000 positive samples from the seed corpus and pair them with 500,000 randomly selected web pages from Common Crawl as negative samples. Training is conducted with the help of an open-source implementation\footnote{\url{https://fasttext.cc}}, where we set the vector dimension to 256, the learning rate at 0.1, a maximum word n-gram length of 3, a minimum word frequency of 3, and 3 training epochs. In order to downsize the original Common Crawl dataset, we apply both URL-based deduplication and near-deduplication techniques, arriving at a set of 40B unique HTML web pages. Using the trained fastText model, we identify potential mathematical pages within this deduplicated subset. To ensure quality, these pages are sorted according to their fastText prediction scores, with only the top-ranking results retained. We determine the optimal volume of data by running pre-training trials on samples containing the top 40B, 80B, 120B, and 160B tokens; in this initial phase, we opt to keep the top 40B tokens.

Following the initial round of data gathering, a substantial number of mathematical web pages are still missing, primarily because the set of positive examples used to train the fastText model does not exhibit enough diversity. To address this limitation, we expand the seed corpus by incorporating additional mathematical web sources, aiming to improve the fastText model’s performance. Concretely, we segment the complete Common Crawl dataset into distinct domains, where a domain comprises web pages under the same base URL. We then determine, for each domain, the proportion of web pages that were collected during the first round. Domains in which more than 10\% of their web pages have already been gathered are deemed math-related (for instance, \url{mathoverflow.net}). 
We proceed to manually label the URLs within these domains that are pertinent to mathematical content, such as \url{mathoverflow.net/questions}. Any uncollected web pages linked to these annotated URLs are subsequently included in the seed corpus. This strategy yields a greater variety of positive examples, allowing us to further refine the fastText model so that it can identify a larger fraction of mathematical data in future rounds. By the end of four iterative data collection cycles, we accumulate 35.5M mathematical web pages, comprising 120B tokens. During the fourth iteration, it becomes evident that close to 98\% of the content had already been harvested in the preceding iteration, leading us to discontinue further data collection.

In order to prevent contamination of benchmarks, we adopt the methodology described in \cite{deepseek-coder}, excluding web content that includes questions or answers originating from English math benchmarks like GSM8K \citep{gsm8k} and MATH \citep{MATH}, as well as Chinese benchmarks such as CMATH \citep{wei2023cmath} and AGIEval \citep{agieval}. Our filtering process operates as follows: any portion of text from our mathematics training data is excluded if it contains a 10-gram sequence that exactly matches any substring present in the evaluation benchmark datasets. For benchmark examples that are less than 10 grams in length but no fewer than 3 grams, we apply an exact matching procedure to eliminate any web pages found to be contaminated.

\subsection{Validating the Quality of the DeepSeekMath~Corpus}

We run pre-training experiments to investigate how the DeepSeekMath~Corpus is compared with the recently released math-training corpora:

\begin{itemize}[topsep=0pt]
    \item \textbf{MathPile} \citep{mathpile}: a collection of texts (8.9B tokens) drawn from multiple sources, such as textbooks, Wikipedia, ProofWiki, CommonCrawl, StackExchange, and arXiv; notably, arXiv is the primary contributor, providing more than 85\% of the dataset;
    \item \textbf{OpenWebMath} \citep{openwebmath}: a subset of CommonCrawl curated for mathematical material, amounting to a total of 13.6B tokens;
    \item \textbf{Proof-Pile-2} \citep{llemma}: a dataset for mathematics that integrates OpenWebMath, AlgebraicStack (comprising 10.3B tokens of math-related code), and papers from arXiv (totaling 28.0B tokens). In line with the approach in \cite{llemma}, our experiments with Proof-Pile-2 employ an arXiv:Web:Code token distribution ratio of 2:4:1.
\end{itemize}

\subsubsection{Training Setting}
\label{sec:quality-policy}
Mathematics-oriented training is conducted on a general-purpose language model comprising 1.3 billion parameters, architecturally consistent with the DeepSeek LLM family \citep{deepseek-llm}; this variant is referred to as DeepSeek-LLM 1.3B. Separate models are trained for each distinct mathematical dataset, each for a total of 150 billion tokens. All training utilizes the streamlined HAI-LLM framework \citep{haillm}. Consistent with the methodology employed for DeepSeek LLMs, the AdamW optimizer \citep{adamW} is utilized with hyperparameters $\beta_1=0.9$, $\beta_2=0.95$, and a $\mathrm{weight\_decay}$ of 0.1. A multi-step learning rate schedule is adopted: the learning rate increases to its maximum following a 2,000-step warmup period, reduces to 31.6\% of its peak by 80\% training completion, and tapers further to 10.0\% of the peak by 90\% completion. The peak learning rate is set to 5.3e-4, with a per-step batch size of 4 million tokens and a maximum context window of 4,096 tokens.

\subsubsection{Evaluation Results}

\textbf{The DeepSeekMath~Corpus is of high quality, covers multilingual mathematical content, and is the largest in size.}

\begin{itemize}[topsep=0pt]
    \item \textbf{High-quality}:
    We assess the downstream abilities of models using eight mathematical benchmarks with few-shot chain-of-thought prompting as detailed in \cite{cot}. According to Table \ref{tab:corpora_comparison}, models trained on the DeepSeekMath~Corpus substantially outperform others. Moreover, Figure \ref{fig:corpora_comparisons} illustrates that the DeepSeekMath-trained model surpasses the performance of the model trained on Proof-Pile-2 after processing 50B tokens (equivalent to one epoch over Proof-Pile-2), which suggests that the overall standard of DeepSeekMath~Corpus is superior.
    \item \textbf{Multilingual}:
    The DeepSeekMath~Corpus includes a broad range of languages, with English and Chinese being the most prevalent. Table \ref{tab:corpora_comparison} demonstrates that training with this corpus significantly improves mathematical reasoning abilities in both English and Chinese settings. By comparison, most current mathematical corpora predominantly focus on English and provide limited benefits for Chinese tasks, occasionally resulting in degraded performance in that language.
    \item \textbf{Large-scale}:
    In terms of volume, the DeepSeekMath~Corpus greatly surpasses existing mathematical corpora. As illustrated by Figure \ref{fig:corpora_comparisons}, when DeepSeek-LLM 1.3B is trained using this dataset, the model exhibits a sharper learning trajectory and achieves more durable gains. In contrast, baseline datasets are considerably smaller and must be revisited numerous times during training, causing rapid stagnation in model performance.
\end{itemize}

\subsection{Training and Evaluating DeepSeekMath-Base 7B}

In this section, we present DeepSeekMath-Base 7B, a foundational model designed for robust mathematical reasoning. The initial weights of our model are derived from DeepSeek-Coder-Base-v1.5 7B \citep{deepseek-coder}, and the training process covers 500B tokens. The training data is sourced from multiple domains: 56\% originates from the DeepSeekMath Corpus, 4\% from AlgebraicStack, 10\% from arXiv, 20\% comprises Github code, while the remaining 10\% includes natural language samples from Common Crawl in both English and Chinese. The training regimen primarily aligns with the methodology outlined in Section \ref{sec:quality-policy}, with adjustments such that the peak learning rate is set to 4.2e-4, and each batch processes 10M tokens.

We present an extensive evaluation of DeepSeekMath-Base 7B’s mathematical proficiency, emphasizing its capacity to independently generate complete mathematical solutions, utilize tools to address mathematical tasks, and perform formal theorem proving. Additionally, we broaden our analysis by offering a general overview of the base model's abilities in natural language understanding, reasoning, and programming.

\paragraph{Mathematical Problem Solving with Step-by-Step Reasoning}

We assess the capabilities of DeepSeekMath-Base in mathematical problem-solving by utilizing few-shot chain-of-thought prompting \citep{cot} across eight benchmarks in both English and Chinese. The selected benchmarks span quantitative reasoning datasets such as GSM8K \citep{gsm8k}, MATH \citep{MATH}, and CMATH \citep{wei2023cmath}, as well as multiple-choice assessments like MMLU-STEM \citep{mmlu} and Gaokao-MathQA \citep{agieval}. This range addresses various mathematical domains, from primary to collegiate levels.

As indicated in Table \ref{tab:base_math_cot}, DeepSeekMath-Base 7B outperforms all other open-source base models on each of the eight evaluated benchmarks, including established models such as Mistral 7B \citep{mistral} and the newly introduced Llemma 34B \citep{llemma}, which has been trained on Proof-Pile-2 \citep{llemma} for mathematical reasoning. Of particular note is its performance on the MATH dataset, where DeepSeekMath-Base exceeds other open-source base models by a margin greater than 10\% in absolute terms, and achieves better results than Minerva 540B \citep{minerva}—a proprietary model that is 77 times larger and enhanced via additional training on mathematical corpora based on PaLM \citep{palm}.

\paragraph{Mathematical Problem Solving with Tool Use} We assess program-assisted mathematical problem-solving abilities on GSM8K and MATH via few-shot program-of-thought prompting \citep{pot,pal}. For each task, models are instructed to generate Python code leveraging modules like \textit{math} and \textit{sympy} to perform advanced calculations. The final solution is determined by executing this code and interpreting its output. According to Table \ref{tab:base_math_pot_proof}, DeepSeekMath-Base 7B surpasses the previous best Llemma 34B in performance.

\paragraph{Formal Mathematics}

The automation of formal proofs plays a critical role in bolstering the correctness and dependability of mathematical reasoning, while also improving productivity—a trend that has received mounting interest in recent times. In our study, we assess DeepSeekMath-Base 7B in the context of informal-to-formal proof generation as described in \citep{dsp_proof}, a task involving the production of a formalized proof utilizing an informal problem statement, its formal expression, and an informal proof. To facilitate this evaluation, we employ the miniF2F \citep{minif2f} dataset, recognized as a benchmark for Olympiad-level formal mathematics, prompting the model with a few exemplars to elicit formal proofs in Isabelle for each test case. Consistent with the methodology in \cite{dsp_proof}, proof sketches are produced by the language model before Sledgehammer \citep{sledgehammer}, an established automated prover, is applied to complete any incomplete aspects of the proofs. As presented in Table \ref{tab:base_math_pot_proof}, our results indicate that DeepSeekMath-Base 7B attains commendable performance in the task of proof autoformalization.

\paragraph{Natural Language Understanding, Reasoning, and Code}

We assess the capabilities of the models in natural language understanding using MMLU \citep{mmlu}, reasoning through BBH \citep{bbh}, and programming skills with HumanEval \citep{codex} and MBPP \citep{mbpp}. As detailed in Table \ref{tab:understanding_reasoning_code}, DeepSeekMath-Base 7B achieves marked improvements over its predecessor, DeepSeek-Coder-Base-v1.5 \citep{deepseek-coder}, on both MMLU and BBH, indicating that mathematical training can enhance both linguistic comprehension and reasoning. Furthermore, thanks to continual training that includes code-related tokens, DeepSeekMath-Base 7B is able to sustain the same coding benchmark performance as DeepSeek-Coder-Base-v1.5 on HumanEval and MBPP. In summary, DeepSeekMath-Base 7B delivers considerably better results than the broadly trained Mistral 7B \citep{mistral} across all three benchmarks targeting reasoning and coding.

\section{Supervised Fine-Tuning}

\subsection{SFT Data Curation}

We develop a diverse instruction-tuning dataset for mathematics that includes problems in both English and Chinese, spanning a range of mathematical domains and degrees of difficulty. Each problem is matched with a corresponding solution in various reasoning formats: chain-of-thought (CoT) \citep{cot}, program-of-thought (PoT) \citep{pot,pal}, and formats that incorporate external tools \citep{tora}. Altogether, the dataset contains 776K examples for training.

\begin{itemize}[topsep=0pt]
    \item \textbf{English mathematical datasets}: Our resources include GSM8K and MATH, each supplemented with tool-based solution annotations. In addition, we utilize selected portions of MathInstruct \citep{MathInstruct} and the Lila-OOD \citep{lila} training set, both providing problems resolved through CoT or PoT techniques. This English dataset compilation encompasses a wide array of mathematical disciplines, such as algebra, probability, number theory, calculus, and geometry.
    \item \textbf{Chinese mathematical datasets}: We assemble a collection of Chinese K-12 mathematics questions, encompassing 76 different sub-fields—including linear equations—with corresponding solutions provided in both CoT and tool-supported reasoning styles.
\end{itemize}

\subsection{Training and Evaluating DeepSeekMath-Instruct 7B}

This section presents DeepSeekMath-Instruct 7B, a model developed by applying mathematical instruction tuning to DeepSeekMath-Base. For training, we concatenate randomly sampled examples, ensuring that the combined sequence does not exceed a context length of 4K tokens. The model is trained for a total of 500 steps using a batch size of 256 and a fixed learning rate of $5 \times 10^{-5}$.

To assess mathematical proficiency, we examine model performance on four quantitative reasoning benchmarks, both in English and Chinese, considering scenarios with and without external tools. Our evaluation includes comparisons with the top-performing contemporaneous models:

\begin{itemize}[topsep=0pt]
    \item \textbf{Closed-source models} comprise the following: (1) the GPT family, with GPT-4 \citep{gpt4} and GPT-4 Code Interpreter~\footnote{\url{https://openai.com/blog/chatgpt-plugins\#\#code-interpreter}} representing the most advanced offerings, (2) Gemini Ultra and Pro \citep{gemini}, (3) Inflection-2 \citep{inflection-2}, (4) Grok-1~\footnote{\url{https://x.ai/model-card}}. Additionally, several Chinese tech firms have introduced (5) Baichuan-3~\footnote{\url{https://www.baichuan-ai.com}}, and (6) the newly launched GLM-4~\footnote{\url{https://open.bigmodel.cn/dev/api\#glm-4}}.

    from the GLM family \citep{glm}.

These models are designed for a wide range of applications and most have been subjected to various alignment processes.

\item \textbf{Open-source models} encompass: general-purpose models such as (1) DeepSeek-LLM-Chat 67B \citep{deepseek-llm}, (2) Qwen 72B \citep{qwen}, (3) SeaLLM-v2 7B \citep{seallm}, and (4) ChatGLM3 6B \citep{chatglm3}. There are also specialized models with mathematical capabilities, including (5) InternLM2-Math 20B~\footnote{\url{https://github.com/InternLM/InternLM-Math}}, which extends InternLM2 by further training on mathematical content and subsequent instruction fine-tuning, (6) Math-Shepherd-Mistral 7B, which leverages PPO training \citep{schulman2017proximal} on Mistral 7B \citep{mistral} utilizing a process-based reward model, (7) the WizardMath family \citep{wizardmath}, which enhances Mistral 7B and Llama-2 70B \citep{llama2} for mathematical reasoning by employing evolve-instruct—a form of instruction tuning based on AI-generated instructions—and PPO training, predominantly drawing problems from GSM8K and MATH, (8) MetaMath 70B \citep{metamath}, an Llama-2 70B model fine-tuned on extended GSM8K and MATH datasets, (9) ToRA 34B \cite{tora}, derived from CodeLlama 34B and further tuned for tool-augmented mathematical reasoning, and (10) MAmmoTH 70B \citep{MathInstruct}, which applies instruction tuning to Llama-2 70B using the MathInstruct dataset.

Table \ref{tab:sft_rl_math} presents the results for evaluation without access to external tools, where DeepSeekMath-Instruct 7B exhibits robust step-by-step reasoning abilities. Specifically, on the challenging MATH dataset designed for competition-level problems, our model outperforms all publicly available models and most commercial systems (such as Inflection-2 and Gemini Pro) by a minimum margin of 9\% absolute. This holds true even in comparison with significantly larger models (e.g., Qwen 72B) or those deliberately improved via reinforcement learning tailored to mathematical tasks (e.g., WizardMath-v1.1 7B). Although DeepSeekMath-Instruct achieves performance on par with proprietary Chinese models such as GLM-4 and Baichuan-3 on the MATH benchmark, it still falls short of the capabilities demonstrated by GPT-4 and Gemini Ultra.

When models are permitted to leverage both natural language reasoning and the utilization of programmatic tools to solve problems, DeepSeekMath-Instruct 7B achieves nearly 60\% accuracy on the MATH benchmark, outperforming all previously available open-source alternatives. Across additional benchmarks, our model performs on par with DeepSeek-LLM-Chat 67B, the former leading model, despite having only one-tenth of its parameters.

\subsection{Group Relative Policy Optimization} Recent research demonstrates that reinforcement learning (RL) substantially enhances the mathematical reasoning skills of large language models (LLMs) following the Supervised Fine-Tuning (SFT) phase \citep{wang2023math,wizardmath}. This section presents our novel RL method, Group Relative Policy Optimization (GRPO), which is both efficient and effective.

\subsubsection{From PPO to GRPO} Proximal Policy Optimization (PPO) \citep{schulman2017proximal} represents a leading actor-critic method in reinforcement learning, and has become the predominant approach for the RL-based fine-tuning of large language models \citep{ouyang2022training}. The optimization of LLMs under this paradigm is carried out by maximizing a specific surrogate objective, defined as follows:

\begin{equation}
\footnotesize
    \mathcal{J}_{PPO}(\theta) = \mathbb{E}{[q \sim P(Q), o \sim \pi_{\theta_{old}}(O|q)]} \frac{1}{|o|} \sum_{t=1}^{|o|} \min \left[ \frac{\pi_\theta(o_{t} | q, o_{<t})}{\pi_{\theta_{old}}(o_{t} | q, o_{<t})} A_{t}, \text{clip} \left( \frac{\pi_\theta(o_{t} | q, o_{<t})}{\pi_{\theta_{old}}(o_{t} | q, o_{<t})}, 1 - \varepsilon, 1 + \varepsilon \right)  A_{t} \right] ,
\end{equation}

Here, $\pi_{\theta}$ denotes the present policy model, while $\pi_{\theta_{old}}$ refers to its prior iteration. The variables $q$ and $o$ correspond to questions and their outputs, sampled from the question corpus and according to the actions of the previous policy $\pi_{\theta_{old}}$, respectively. The parameter $\varepsilon$ serves as a clipping hyper-parameter introduced within PPO to promote stability during the learning process. The advantage term $A_t$ is estimated through Generalized Advantage Estimation (GAE) \citep{gae}, utilizing the subsequent rewards $\{r_{\ge t}\}$ and a value function $V_{\psi}$ that is learned concurrently. Consequently, PPO mandates training a value function in tandem with the policy network. To counteract excessive optimization of the reward model, a prevalent method is to augment the per-token reward with a KL divergence penalty relative to a reference model at each token \citep{ouyang2022training}, as follows:

\begin{equation}
    r_{t} = r_\varphi(q, o_{\le t}) - \beta \log\frac{\pi_{\theta}(o_{t}|q, o_{<t})}{\pi_{ref}(o_{t}|q, o_{<t})},
\label{eq:PPO-reward}
\end{equation}

where $r_\varphi$  is the reward model, $\pi_{ref}$ is the reference model, which is usually the initial SFT model, and $\beta$  is the coefficient of the KL penalty.

In standard PPO implementations, the value function is usually represented by a neural network of similar size to the policy model, resulting in significant demands on memory and computation. Furthermore, within RL training, the value function serves as a baseline during advantage computation to decrease variance. However, in the context of LLMs, the reward model often assigns a score solely to the final token, posing additional challenges for developing a value function that accurately predicts at every token position. To overcome these limitations, we introduce Group Relative Policy Optimization (GRPO), as depicted in Figure \ref{fig:grpo}. GRPO eliminates the necessity for a separate value function as required in PPO. Instead, it defines the baseline as the mean reward obtained from multiple sampled responses to the same prompt. Specifically, given a question $q$, GRPO generates a collection of outputs $\{o_1, o_2, \cdots, o_G\}$ using the previous policy $\pi_{\theta_{old}}$ and trains the policy model to optimize the following objective:

\begin{equation}
\footnotesize
\begin{split}
    \mathcal{J}_{GRPO}(\theta) &= \mathbb{E}{[q \sim P(Q), \{o_i\}_{i=1}^G \sim \pi_{\theta_{old}}(O|q)]}  \\
    & \frac{1}{G}\sum_{i=1}^G\frac{1}{|o_i|} \sum_{t=1}^{|o_i|} \Bigg\{ \min \left[ \frac{\pi_\theta(o_{i,t} | q, o_{i,<t})}{\pi_{\theta_{old}}(o_{i,t} | q, o_{i,<t})} \hat{A}_{i,t},\, \text{clip} \left( \frac{\pi_\theta(o_{i,t} | q, o_{i,<t})}{\pi_{\theta_{old}}(o_{i,t} | q, o_{i,<t})}, 1 - \varepsilon, 1 + \varepsilon \right)  \hat{A}_{i,t} \right] \\
    & \hspace{6cm} - \beta \mathbb{D}_{KL}\left[\pi_{\theta} || \pi_{ref}\right]\Bigg\} ,
\end{split}
\label{eq:GRPO-obj}
\end{equation}

where $\varepsilon$ and $\beta$ denote hyper-parameters, and $\hat{A}_{i,t}$ refers to the advantage, which is computed solely from the relative rewards among the outputs within each group. This approach will be elaborated upon in the subsequent subsections. By using within-group comparisons to estimate advantages, GRPO captures the inherently comparative objective of reward models, which are commonly trained using datasets featuring output pairwise comparisons for the same prompt. Additionally, GRPO differs from prior methods that incorporate KL penalties directly into the reward function; instead, it enforces regularization by explicitly including the KL divergence between the current policy and the reference policy within the loss term. This strategy streamlines the computation of $\hat{A}_{i,t}$, since the KL does not interact with its estimation. Unlike the KL penalty present in (\ref{eq:PPO-reward}), here the KL divergence is computed with the following unbiased estimator \citep{kl_approx}:

\begin{equation}
\small
    \mathbb{D}_{KL}\left[\pi_{\theta} || \pi_{ref}\right] = \frac{\pi_{ref}(o_{i,t}|q,o_{i,<t})}{\pi_{\theta}(o_{i,t}|q,o_{i,<t})}- \log\frac{\pi_{ref}(o_{i,t}|q,o_{i,<t})}{\pi_{\theta}(o_{i,t}|q,o_{i,<t})} - 1,
\end{equation}

which is guaranteed to be positive.

\begin{algorithm}[t]
  \small
  \caption{Iterative Group Relative Policy Optimization}
  \textbf{Input} initial policy model $\pi_{\theta_{\text{init}}}$; reward models $r_\varphi$; task prompts $\mathcal{D}$; 
  hyperparameters $\varepsilon$, $\beta$, $\mu$
  \begin{algorithmic}[1]
    \State Initialize policy model: $\pi_\theta \leftarrow \pi_{\theta_{\text{init}}}$
    \For{iteration = 1, \dots, I}
       \State Assign the current policy model as the reference: $\pi_{ref} \leftarrow \pi_{\theta}$
      \For{step = 1, \dots, M}
      \State Randomly select a batch $\mathcal{D}_b$ from $\mathcal{D}$
      \State Assign latest policy to old policy: $\pi_{\theta_{old}} \leftarrow \pi_{\theta}$
      \State For each query $q \in \mathcal{D}_b$, generate $G$ responses $\{o_i\}_{i=1}^G \sim \pi_{\theta_{old}} (\cdot \mid q)$
      \State Evaluate each sampled response $o_i$ to obtain rewards $\{r_i\}_{i=1}^{G}$ using $r_{\varphi}$ 
      \State For every response $o_i$, estimate the group relative advantage $\hat{A}_{i,t}$ at every token position $t$
      \For{GRPO iteration = 1, \dots, $\mu$}
        \State Perform policy update by optimizing the GRPO loss, as shown in Equation~\ref{eq:GC-GRPO}
      \EndFor
    \EndFor 
    \State Enhance reward model $r_\varphi$ via ongoing training leveraging a replay buffer.
    \EndFor 
  \end{algorithmic}
  \textbf{Output} $\pi_\theta$
  \label{alg:iter-grpo}
\end{algorithm}

\subsubsection{Outcome Supervision RL with GRPO} Specifically, given a question $q$, a collection of outputs $\{o_1, o_2, \cdots, o_G\}$ is generated using the previous policy model $\pi_{\theta_{old}}$. Each of these outputs is evaluated using a reward model, resulting in a set of $G$ rewards $\mathbf{r}=\{r_1, r_2, \cdots, r_G\}$. The obtained rewards are normalized by first removing the mean of the group and then dividing by the standard deviation. In outcome supervision, this normalized reward is assigned at the conclusion of each output $o_i$, with every token in $o_i$ sharing the same advantage value equal to the normalized reward; that is, $\hat{A}_{i, t} = \widetilde{r}_i = \frac{r_i- {\rm mean}(\mathbf{r})}{{\rm std}(\mathbf{r})}$. The policy is then updated to maximize the objective provided in equation (\ref{eq:GRPO-obj}).

\subsubsection{Process Supervision RL with GRPO} 

While outcome supervision delivers a reward solely at the conclusion of an output sequence, this sparse feedback is often inadequate or inefficient for guiding policy learning on intricate mathematical problems. Building on the approach proposed by \cite{wang2023math}, we additionally consider process supervision, where feedback is provided at the completion of every reasoning step. Formally, for a given question $q$ and a collection of $G$ sampled outputs $\{o_1, o_2, \cdots, o_G\}$, a process reward model assigns a score to each intermediate step of each output, producing the set of rewards: $\mathbf{R} = \{ \{r_1^{index(1)},\cdots,r_1^{index(K_1)}\}, \cdots,  \{r_G^{index(1)},\cdots,r_G^{index(K_G)}\} \}$, where $index(j)$ corresponds to the token marking the end of the $j$-th step, and $K_i$ denotes the number of reasoning steps in the $i$-th generated output. To ensure comparability, these rewards are standardized using the mean and standard deviation as follows: $\widetilde{r}_i^{index(j)} = \frac{r_i^{index(j)} - {\rm mean(\mathbf{R})}}{{\rm std(\mathbf{R})}}$.

The process supervision framework then defines the advantage for each token as the cumulative sum of the normalized rewards assigned to all subsequent reasoning steps, that is, $\hat{A}_{i, t} = \sum_{index(j) \ge t} \widetilde{r}_i^{index(j)}$. The resulting advantage estimates are employed to update the policy by maximizing the objective presented in equation (\ref{eq:GRPO-obj}).

\subsubsection{Iterative RL with GRPO} During the reinforcement learning process, the reward model initially used may become inadequate for effectively guiding the latest policy model. To address this, we investigate iterative RL with GRPO. As illustrated in Algorithm \ref{alg:iter-grpo}, this approach constructs updated training datasets for the reward model by utilizing samples drawn from the current policy model. The reward model is then further optimized through continual training with a replay strategy, where 10\% of previously collected data is maintained in the dataset. Afterward, the policy model is set as the new reference model, and training proceeds by updating the policy model in accordance with the newly improved reward model.

\subsection{Training and Evaluating DeepSeekMath-RL}

Our reinforcement learning experiments utilize DeepSeekMath-Instruct 7B as the policy backbone. The RL dataset is composed of approximately 144K chain-of-thought-format questions extracted from the SFT data pertaining to GSM8K and MATH, while all other SFT questions are deliberately omitted to analyze the effects of RL on benchmarks without data exposure during training. The methodology for assembling the reward model training dataset adheres to \citep{wang2023math}. The reward model is initially trained on DeepSeekMath-Base 7B, employing a learning rate of 2e-5. For GRPO, the policy model operates with a learning rate of 1e-6, a KL coefficient of 0.04, and, per question, we generate $64$ response samples. The maximum sequence length permitted is 1024 tokens, with a corresponding batch size of 1024. Policy parameters are updated just once after every exploration step. DeepSeekMath-RL 7B undergoes evaluation on the same benchmarks as DeepSeekMath-Instruct 7B. For DeepSeekMath-RL 7B, tasks involving GSM8K and MATH chain-of-thought reasoning are classified as in-domain, while the rest of the benchmarks are considered out-of-domain.

Table \ref{tab:sft_rl_math} presents comparative results for open- and closed-source models that utilize chain-of-thought and tool-augmented reasoning on both English and Chinese evaluation datasets. Our analysis indicates the following: 1) DeepSeekMath-RL 7B achieves accuracy rates of 88.2\% on GSM8K and 51.7\% on MATH when employing chain-of-thought reasoning, outperforming every open-source model within the 7B–70B parameter range, and exceeding most closed-source competitors as well. 2) Notably, the training for DeepSeekMath-RL 7B exclusively uses chain-of-thought formatted instruction tuning data from GSM8K and MATH, building upon the foundation of DeepSeekMath-Instruct 7B. Despite its narrower training data, DeepSeekMath-RL 7B consistently achieves higher scores than DeepSeekMath-Instruct 7B in all evaluation categories, underscoring the advantages conferred by reinforcement learning.

\section{Discussion}
In this section, we will share our findings in pre-training and RL experiments. 

\subsection{Lessons Learnt in Pre-Training}

We begin by detailing our observations from the pre-training phase. Unless noted differently, all training procedures follow the configurations described in Section \ref{sec:quality-policy}. Notably, references to the DeepSeekMath Corpus within this section pertain to a dataset consisting of 89B tokens gathered during the second round of data collection.

\subsubsection{Code Training Benefits Mathematical Reasoning}

An often-cited but unproven hypothesis posits that practicing code enhances reasoning skills. In this study, we seek to provide partial evidence regarding this claim, focusing specifically on mathematics: engaging in code training enhances models’ proficiency in mathematical reasoning, regardless of whether auxiliary tools are used.

To study how code training affects mathematical reasoning, we experimented with the following two-stage training and one-stage training settings:

\textbf{Two-Stage Training}

\begin{itemize}[topsep=0pt]
    \item \textbf{Code Training for 400B Tokens $\rightarrow$ Math Training for 150B Tokens}: The DeepSeek-LLM 1.3B model undergoes a pretraining phase with 400B code tokens, which is then succeeded by a subsequent stage involving 150B math tokens.
    \item \textbf{General Training for 400B Tokens $\rightarrow$ Math Training for 150B Tokens}: For comparative analysis, we conduct a control experiment wherein, during the initial training stage, the model is exposed to 400B general tokens (sourced from an extensive, large-scale corpus curated by DeepSeek-AI), rather than code tokens. This setup is designed to assess whether pretraining on code provides any specific benefits for enhancing mathematical reasoning capabilities relative to a generic dataset.
\end{itemize}

\textbf{One-Stage Training}

\begin{itemize}[topsep=0pt]
    \item \textbf{Math Pretraining over 150B Tokens}: DeepSeek-LLM 1.3B undergoes pretraining using 150B tokens from mathematical data;
    \item \textbf{Integrated Training with 400B Code Tokens and 150B Math Tokens}: Sequentially pretraining on math data after code data diminishes coding ability. Therefore, we assess whether combining math and code tokens within a unified training phase enhances mathematical reasoning while mitigating catastrophic forgetting observed with multi-stage training.
\end{itemize}

\paragraph{Results}
Table \ref{tab:code-math} and Table \ref{tab:code-math-general-eval} demonstrate the downstream performance under different training settings.

Code training facilitates program-aided mathematical reasoning across both one-stage and two-stage training paradigms. As demonstrated in Table \ref{tab:code-math}, applying code training during the two-stage approach already results in marked improvements in solving GSM8K and MATH tasks with Python. Introducing a subsequent math training phase leads to additional gains. Notably, in the one-stage setting, jointly presenting code and math tokens helps prevent the catastrophic forgetting typically observed in two-stage training and further boosts performance in both coding tasks (refer to Table \ref{tab:code-math-general-eval}) and program-aided mathematical reasoning tasks (see Table \ref{tab:code-math}).

Training with code also enhances mathematical reasoning abilities even in the absence of tool use. In the two-stage training paradigm, the preliminary phase focused on code already produces noticeable improvements. Moreover, it facilitates more effective learning during the following math training stage, which ultimately yields superior results. In contrast, merging code tokens and math tokens within a single training stage diminishes mathematical reasoning performance when tools are not used. A possible explanation is that DeepSeek-LLM 1.3B, constrained by its relatively small model size, may not possess sufficient capacity to effectively learn from both code and mathematics data at the same time.

\subsubsection{ArXiv Papers Seem Ineffective in Improving Mathematical Reasoning}

ArXiv papers are frequently utilized as part of mathematical pre-training datasets \citep{minerva,gpt-f,llemma,mathpile}. Nonetheless, there is a lack of comprehensive investigation into how these papers influence mathematical reasoning. Somewhat unexpectedly, our empirical results suggest that including arXiv papers yields little to no benefit for mathematical reasoning capabilities. To explore this further, we employ models of various scales—including DeepSeek-LLM 1.3B and DeepSeek-Coder-Base-v1.5 7B \citep{deepseek-coder}—and assess their performance on arXiv text subjected to different preprocessing workflows:

\begin{itemize}[topsep=0pt]
    \item \textbf{MathPile} \citep{mathpile}: a dataset consisting of 8.9B tokens, constructed through heuristic-based cleaning and filtering procedures, with scientific articles from arXiv comprising more than 85\% of its content;
    \item \textbf{ArXiv-RedPajama} \citep{redpajama}: encompasses all arXiv LaTeX files, carefully stripped of preambles, comments, macros, and bibliographic entries, yielding a total of 28.0B tokens.
\end{itemize}

During our experiments, we individually trained DeepSeek-LLM 1.3B with 150B tokens and DeepSeek-Coder-Base-v1.5 7B with 40B tokens on the respective arXiv datasets. Our findings indicate that arXiv papers provide little to no benefit for enhancing mathematical reasoning abilities. When models were trained exclusively on the arXiv corpus, neither demonstrated significant gains— and in some cases, their performance on a range of mathematical benchmarks declined. The benchmarks assessed span diverse areas of complexity, including quantitative reasoning tasks such as GSM8K and MATH (Table \ref{tab:arxiv-cot}), multiple-choice tests exemplified by MMLU-STEM (Table \ref{tab:arxiv-cot}), and formal mathematics evaluations such as miniF2F (Table \ref{tab:arxiv_atp}).

However, this conclusion has its limitations and should be taken with a grain of salt. We have not yet studied:

\begin{itemize}[topsep=0pt]
    \item How arXiv tokens influence math-specific tasks beyond the scope of this study, for example, the informalization of theorems, which involves translating formal statements or proofs into informal language;
    \item The influence of integrating arXiv tokens with additional data modalities;
    \item If scaling up the model size amplifies the advantages provided by incorporating arXiv papers.
\end{itemize}

Thus, further exploration is required, which we leave for future studies.

\subsection{Insights of Reinforcement Learning}
\subsubsection{Towards to a Unified Paradigm} In this section, we introduce a cohesive framework to examine various training strategies, encompassing SFT, RFT, DPO, PPO, and GRPO. We also empirically investigate the components that constitute this unified paradigm. In essence, for many training algorithms, the gradient of the objective with respect to parameters $\theta$ can be formalized as follows:

\begin{equation}
    \nabla_{\theta}\mathcal{J}_{\textcolor{red}{\mathcal{A}}}(\theta) = \mathbb{E}[\underbrace{(q,o) \sim \textcolor{red}{\mathcal{D}}}_{Data \ Source}]\left( \frac{1}{|o|} \sum_{t=1}^{|o|}  \underbrace{GC_{{\mathcal{A}}}(q, o, t, \textcolor{red}{\pi_{{rf}}})}_{Gradient \ Coefficient}  \nabla_{\theta}\log \pi_{\theta}(o_t | q, o_{<t})\right).
\label{eq:objective}
\end{equation}

Three essential elements are involved: 1) \textit{Data Source $\mathcal{D}$}, specifying the dataset used for training; 2) \textit{Reward Function $\pi_{{rf}}$}, providing the signal used to assign rewards during training; and 3) \textit{Algorithm $\mathcal{A}$}, which utilizes both the data and the reward signal to compute the gradient coefficient $GC$, thereby establishing the level of penalization or reinforcement applied to the data. Using this unified framework, we examine a range of representative approaches:

\begin{itemize}
    \item \textbf{Supervised Fine-tuning (SFT)}: In SFT, a pretrained model is further trained on a dataset curated by humans specifically for SFT.
    \item \textbf{Rejection Sampling Fine-tuning (RFT)}: RFT builds on the SFT model by conducting an additional fine-tuning phase, where only selected responses generated by the SFT model in response to SFT prompts—those deemed correct—are used for training.
    \item \textbf{Direct Preference Optimization (DPO)}: DPO continues the refinement of the SFT model, leveraging an expanded set of outputs produced by the SFT model and employing a pairwise DPO loss during fine-tuning.
    \item \textbf{Online Rejection Sampling Fine-tuning (Online RFT)}: Unlike standard RFT, Online RFT starts with the SFT model as the policy and updates it iteratively, fine-tuning on newly generated responses from the current version of the policy model in real time.
    \item \textbf{PPO/GRPO}: With PPO/GRPO, the policy model begins as a copy of the SFT model and then undergoes reinforcement learning, using outputs sampled directly from the up-to-date policy.
\end{itemize}

We summarize the components of these methods in Table \ref{tab:data_policy}.
Please refer to Appendix \ref{app:analysis-rl} for a more detailed derivation process.

\paragraph{Observation about Data Source} The data source is categorized into two types: online sampling and offline sampling. In the online sampling setup, training data is derived directly from trajectories generated by the policy model as it trains in real time. Conversely, offline sampling refers to training data obtained from samples produced by the pretrained SFT model. RFT and DPO utilize offline sampling approaches, whereas Online RFT and GRPO make use of online sampling methodologies.

As illustrated in Figure \ref{fig:rl-analysis}, our results indicate that Online RFT substantially surpasses RFT on both benchmark tasks. In particular, although Online RFT performs on par with RFT during the early phases of training, it achieves a pronounced lead in subsequent stages, highlighting the benefits of the online training paradigm. This observation is expected: initially, both the actor and the SFT model behave similarly, and the distinctions captured in sampled data are minimal. As training proceeds, the divergence in data sampled from the actor grows, and the advantages of real-time data collection become increasingly pronounced.

\paragraph{Observation about Gradient Coefficient} In the update step for model parameters, the algorithm utilizes the processed reward signal as a gradient coefficient. Within our experimental setup, we decompose the reward function into two components: `Rule' and `Model'. The `Rule' component evaluates the quality of responses through answer correctness, whereas the `Model' component involves the training of a reward model tasked with scoring individual responses. This reward model is supervised using data generated from the `Rule' judgments. As highlighted in Equations \ref{eq:GC-RFT} and \ref{eq:GC-GRPO}, a fundamental distinction arises between GRPO and Online RFT: GRPO explicitly calibrates its gradient coefficient in accordance with the output of the reward model. This enables it to selectively amplify or diminish feedback for responses, proportionate to the assessed reward magnitude. On the other hand, Online RFT does not possess this adjustment; it refrains from penalizing incorrect answers and applies a consistent reinforcement to all correctly-answered responses, irrespective of their individual differences.

Figure \ref{fig:rl-analysis} illustrates that GRPO outperforms online RFT, underscoring the effectiveness of modifying the positive and negative gradient coefficients. Moreover, GRPO+PS achieves better outcomes than GRPO+OS, demonstrating the advantage of employing finely-tuned, step-sensitive gradient coefficients. Additionally, we investigate iterative RL by performing two rounds of updates in our experiments. The results, depicted in Figure \ref{fig:iter-rl}, reveal that iterative RL leads to considerable performance gains, with the most pronounced improvement observed after the initial iteration.

\subsubsection{Why RL Works?} In this study, we apply reinforcement learning using a selected portion of instruction tuning data and observe notable improvements compared to the baseline instruction-tuned model. To better elucidate the mechanism behind the effectiveness of reinforcement learning, we assess Pass@K and Maj@K accuracy metrics for both Instruct and RL models across two benchmark datasets. Figure \ref{fig:maj-pass} demonstrates that while RL substantially improves Maj@K, it does not affect Pass@K. This suggests that reinforcement learning increases the robustness of the output distribution, \textbf{implying that the gains are mainly due to improved selection of correct answers among TopK candidates, rather than a fundamental enhancement of core modeling abilities.} A similar observation is reported in \citep{wang2023making}, who describe a \textbf{misalignment problem} in reasoning-oriented tasks with SFT models; their results indicate that preference alignment techniques can effectively improve the reasoning capacities of SFT models \citep{yuan2023rrhf,song2023preference,wang2023making}.

\subsubsection{How to Achieve More Effective RL?}
\label{sec:effec_rl}
Our findings illustrate that RL can be highly effective for tasks involving mathematical reasoning. Furthermore, we introduce a consolidated framework that facilitates a coherent understanding of several prominent training methodologies. According to this framework, these approaches can be categorized as either explicit or streamlined versions of RL. As outlined in Equation \ref{eq:objective}, three foundational elements are identified: Data Source, Algorithm, and Reward Function. We also outline possible avenues for advancing each of these components in future work.

\paragraph{Data Source} The data source forms the foundational input for all training frameworks. Within the scope of RL, we define the data source as comprising unlabeled question prompts, along with outputs generated from the policy model via sampling. For our experiments, we restrict ourselves to using questions stemming from the instruction tuning phase and employ simple nucleus sampling when generating responses. We posit that this may underlie the RL pipeline’s limitation in enhancing only Maj@K scores. Future work aims to investigate the behavior of our RL approach on novel, out-of-distribution prompts, alongside the integration of \textbf{more sophisticated sampling (decoding) strategies}, such as those informed by tree-search techniques \citep{yao2023tree}. Furthermore, advances in \textbf{efficient inference techniques} \citep{xia-etal-2023-speculative,leviathan2023fast,kwon2023efficient,xia2024unlocking}, which are crucial for maximizing the exploration capability of policy models, are expected to significantly influence results as well.

\paragraph{Algorithms} Algorithms utilize both the input data and the reward signals to compute gradient coefficients that facilitate updating model parameters. As derived from Equation \ref{eq:objective}, contemporary methods generally exhibit complete \textbf{TRUST} in the feedback provided by the reward function, adjusting the conditional likelihood of particular tokens accordingly. Nevertheless, the assumption of consistently reliable reward signals does not hold, particularly in highly intricate tasks. For instance, even rigorously curated datasets such as PRM800K \citep{lightman2023let}, created by skilled annotators, still include nearly 20\% erroneous labels\footnote{\url{https://github.com/openai/prm800k/issues/12\#issuecomment-1728491852}}. In light of this, our focus turns to reinforcement learning algorithms engineered to be resilient in the presence of noisy reward feedback. We anticipate that \textbf{WEAK-TO-STRONG} alignment strategies \citep{burns2023weak} have the potential to significantly transform the landscape of learning algorithms.

\paragraph{Reward Function} The reward function serves as the primary source of training signals in reinforcement learning. In most RL scenarios, the reward function is represented by a neural reward model. We identify three critical avenues for the advancement of reward models: 1) \textbf{Enhancing the generalization capability of reward models.} It is essential for reward models to generalize beyond seen data, accommodating out-of-distribution tasks and more sophisticated decoding outputs; failure to do so risks that reinforcement learning only reinforces the pre-existing output distribution of LLMs without promoting substantive progress in capability. 2) \textbf{Representing and leveraging reward model uncertainty.} By capturing uncertainty, reward models may serve as an intermediary between weaker reward estimates and the progression toward more robust learning frameworks; 3) \textbf{Scalable development of high-quality process reward models,} which are necessary for delivering detailed and actionable training signals throughout the reasoning steps of complex tasks \citep{lightman2023let,wang2023math}.

\section{Conclusion, Limitation, and Future Work}
We introduce DeepSeekMath, a model that surpasses all other open-source alternatives on the competition-grade MATH benchmark and nearly matches the achievements of proprietary models. The model is based on DeepSeek-Coder-v1.5 7B, which is continually trained over 500B tokens, incorporating a substantial fraction—120B tokens—of mathematical content gathered from Common Crawl. Through comprehensive ablation analysis, we observe that online web pages hold considerable promise as a source of high-quality mathematical materials, whereas data from arXiv do not yield as much benefit as initially anticipated. Our work proposes Group Relative Policy Optimization (GRPO), an adaptation of Proximal Policy Optimization (PPO) that boosts mathematical reasoning performance while lowering memory demands. Experimental findings indicate that GRPO remains beneficial even when DeepSeekMath-Instruct 7B attains high benchmark results. Additionally, we propose a unified conceptual framework for analyzing a range of methods, and we outline various avenues to pursue more effective reinforcement learning strategies.

While DeepSeekMath~demonstrates strong performance on benchmarks assessing quantitative reasoning, its abilities in geometry and theorem-proving lag behind those of closed-source models. For example, our preliminary experiments revealed that DeepSeekMath~struggles with problems involving triangles and ellipses, suggesting a potential data selection bias during both pre-training and fine-tuning phases. Additionally, owing to the limited scale of the model, DeepSeekMath~underperforms relative to GPT-4 when it comes to few-shot learning; whereas GPT-4 benefits from few-shot examples, DeepSeekMath~exhibits little difference between its zero-shot and few-shot results. Moving forward, we aim to enhance our data selection mechanisms to create a higher-quality pre-training corpus, and will investigate approaches discussed in Section \ref{sec:effec_rl} to further advance the reinforcement learning processes for LLMs.

\end{document}